\thispagestyle{plain}

% Single-spaced with tighter paragraph separation. These MUST stay inside
% \begingroup...\endgroup: \include issues a \clearpage but does not open a
% group, so without this the settings leak into the rest of the thesis and
% silently reflow every following chapter.
\begingroup
\singlespacing
\setlength{\parskip}{2pt}

\begin{center}
    \Large
    \thesistitle

    \vspace{0.2cm}
    \large
    \thesisauthorname

    \vspace{0.2cm}
    \large
    Master of Sciences Thesis

    \vspace{0.1cm}
    \large
    Ben-Gurion University of the Negev

    \vspace{0.1cm}
    \large
    \thesisyear

    \vspace{0.3cm}
    \Large
    \textbf{Abstract}
\end{center}

This thesis asks whether short-horizon complication prediction on hospitalized patients can be shaped as temporal next-token prediction over a stream of clinically named intervals, using a transformer-based model that reasons over acquired medical knowledge --- and, as a secondary goal, whether adherence to clinical guidelines governing the same care can be assessed automatically at national scale. Complications such as those governed by inpatient diabetes guidelines depend on the timely arrival of clinical data, the medical knowledge that turns those data into named concepts, and the adherence to care recommendations that connect the two. Predicting these complications, and assessing that adherence and its effect on the same complications, are the two challenges this work takes on. Both are evaluated on hospitalized patients with a diagnosis of diabetes or documented inpatient hyperglycemia. Inpatient hyperglycemia, hypoglycemia, renal failure, and other complications, including death, occur commonly in such patients and are harmful; their management is governed by clinical guidelines whose adherence is rarely measured, which is what makes automated assessment worthwhile in its own right.

A single representation addresses both objectives. A diabetes knowledge base created for the current study encodes the clinical semantics of inpatient care and the Israeli national position paper as quality-assessment temporal patterns. A temporal-abstraction engine converts an OMOP-formatted record into knowledge-based temporal abstractions and, in the same pass, scores how well the recorded care adhered to the guideline's specifications, on a scale from $0$ (lowest compliance) to $1$ (highest compliance). Every interval becomes a named clinical token, so prediction and measurement read one stream.

A \emph{methodological study} on $57{,}078$ MIMIC-IV admissions established the architecture. \textbf{INTERVenE}, a new transformer family operating natively on those intervals, comprises a bidirectional variant for single-pass prediction and an autoregressive variant that generates abstraction trajectories with per-step risk curves. Across the six target complications --- death, kidney complication, hyperglycemia, severe hyperglycemia, hypoglycemia, and severe hypoglycemia --- the encoder reached a support-weighted AUPRC of $0.672$, improving on the strongest neural baseline by $0.041$ with non-overlapping confidence intervals and on a symbolic pipeline consuming the identical interval stream by $0.085$. Sparse-transcoder attribution, computed on the deployed model through a gradient-decoupled hook, resolved every explanation to a named clinical concept. Ablations separated the sources of that gain: collapsing every interval to a point event degraded performance sharply, while replacing the clinically named intervals with an automatically derived $\sigma$-bin discretization did not. On this cohort, knowledge-based abstraction is therefore an interpretability substrate obtained at essentially no cost in accuracy rather than a predictive prerequisite.

A \emph{clinical study} on $48{,}354$ admissions across six Israeli general hospitals carried the same machinery into routine practice. Adherence levels differed by the \emph{kind} of action a recommendation required rather than by its clinical weight: recommendations satisfied by a single discrete action on admission were near-universally implemented, while those calling for repeated insulin decisions across a stay scored as low as $0.15$ and were most often entirely absent. Adherence rose modestly after implementation of a national quality measure, with marked inter-hospital heterogeneity that these data do not explain. Under propensity weighting, above-median adherence was associated with a $44\%$ lower probability of $30$-day mortality and with an $82\%$ lower probability of hyperglycemia, against a $20\%$ rise in mild hypoglycemia and no detectable rise in severe hypoglycemia --- the lopsided trade-off these recommendations are designed to strike. The recommendations that were least adhered to proved to be among those most strongly associated with outcomes.

A final experiment integrated the two studies described above. INTERVenE-Enc transferred to the clinical cohort, reaching an AUPRC of $0.527$ against $0.472$ for the strongest baseline. Here the knowledge base did carry predictive weight: the same $\sigma$-bin substitution cost $0.073$ AUPRC and left that arm level with the raw-stream baselines rather than ahead of them, so the MIMIC-IV equivalence bounds what clinical knowledge contributes when ported to a new source rather than what it contributes generally. A seven-arm ablation ran each input representation with and without the measured adherence profile supplied to the model as input. Supplying it did not shift the aggregate metrics measurably, though the interval-token arm gained $1$--$4$\,pp AUPRC on the six complications the guideline actually targets, while LogReg and set-attention baselines showed no measurable QA effect. Attribution confirmed the model reads the compliance patterns and ranks them among the top drivers for seven of eight outcomes. Adherence is largely noted implicitly through the input intervals the model had processed, and slight residual lift materializes only where attention can weight compliance patterns per outcome --- a statement about representation, not about whether following the guideline matters clinically, which the propensity-weighted analysis addressed separately, finding higher adherence consistently associated with lower mortality and reduced hyperglycemic complications after balancing measured baseline characteristics.

\endgroup
